\documentclass[11pt]{article}
\usepackage[margin=0.85in]{geometry}
\usepackage[utf8]{inputenc}
\usepackage[T1]{fontenc}
\usepackage{amsmath, amssymb}
\usepackage{enumitem}
\usepackage{microtype}
\usepackage{hyperref}

\setlength{\parskip}{4pt}
\setlength{\parindent}{0pt}

\title{\vspace{-2.5em}\textbf{Project Proposal: Tree-HyperLISTA --- Ultra-Lightweight Deep Unfolding for Tree-Sparse Signal Recovery}\vspace{-1em}}
\author{}
\date{}

\begin{document}
\maketitle
\vspace{-2.5em}

\textbf{Group Members:}
\begin{itemize}[nosep, leftmargin=1.5em]
    \item Shresth Verma \hfill Roll No: 22B2211
    \item Aditya Agrawal \hfill Roll No: 22B0955
\end{itemize}

\textbf{Paper to Implement:}\\
\textit{``Hyperparameter Tuning is All You Need for LISTA''} by Xiaohan Chen, Jialin Liu, Zhangyang Wang, and Wotao Yin. Published at \textbf{NeurIPS 2021}. [\href{https://arxiv.org/abs/2110.15900}{arXiv:2110.15900}]

This paper proposes HyperLISTA, a deep unfolding network for sparse signal recovery ($\mathbf{y} = \mathbf{A}\mathbf{x} + \boldsymbol{\varepsilon}$) that reduces the learnable parameters of LISTA (millions) to just \textbf{3 global hyperparameters} $(c_1, c_2, c_3)$, from which all per-layer step sizes, thresholds, and momentum terms are derived via adaptive formulas. It achieves performance comparable to fully-learned LISTA while being far more robust to distribution shift. However, HyperLISTA uses elementwise soft-thresholding and is unaware of \emph{structured} sparsity patterns.

\textbf{Proposed Extension:}\\
We will extend HyperLISTA to \textbf{tree-structured sparsity}, where the signal's nonzero support must form a connected rooted subtree (e.g., wavelet coefficients of images). Our method, \textbf{Tree-HyperLISTA}, will replace the elementwise proximal operator with a tree-consistent operator comprising: (i)~bottom-up subtree scoring, (ii)~greedy tree-consistent support selection with ancestor closure, and (iii)~soft-thresholding within the selected support---while retaining only 3 hyperparameters. We will also implement baselines: Tree-ISTA/FISTA, Tree-IHT, Model-Based CoSaMP (Baraniuk et al., 2010), LISTA, and Tree-LISTA.

\textbf{Datasets:}
\begin{itemize}[nosep, leftmargin=1.5em]
    \item \textbf{Synthetic tree-sparse data}: Balanced binary trees (depth 7--10, $n = 255$ to $1023$), Gaussian sensing matrices, $K$-tree-sparse signals at various SNR levels. Multiple random seeds for statistical robustness.
    \item \textbf{Wavelet-domain image compressed sensing}: Standard Set11 test images, decomposed via 2D Haar wavelet transform. Wavelet coefficients naturally exhibit tree sparsity (coarse-to-fine parent-child structure). CS ratios: 10\%, 25\%, 50\%.
\end{itemize}

\textbf{Evaluation and Validation Strategy:}
\begin{itemize}[nosep, leftmargin=1.5em]
    \item \textbf{Reconstruction quality}: Normalized Mean Squared Error (NMSE in dB) for synthetic data; PSNR and SSIM for image experiments.
    \item \textbf{Support recovery}: Node-level precision and recall to measure whether the correct tree support is identified.
    \item \textbf{Convergence analysis}: NMSE vs.\ layer/iteration curves to verify that Tree-HyperLISTA converges faster than classical baselines and comparably to heavy-parameterized learned methods.
    \item \textbf{Robustness / mismatch testing}: Evaluate all models on test distributions that differ from training (different SNR, sparsity level, perturbed sensing matrix) to validate generalization---a key claimed advantage of HyperLISTA-style methods.
    \item \textbf{Ablation studies}: Sensitivity to each hyperparameter $c_1, c_2, c_3$; comparison of tree support selection mechanisms; effect of descendant decay weight $\rho$.
    \item \textbf{Parameter efficiency}: Compare number of trainable parameters across all methods to demonstrate that 3 parameters with the right structural bias suffice.
\end{itemize}

\end{document}
